\documentclass[a4paper]{article}
\usepackage[utf8]{inputenc}
\usepackage[T1]{fontenc}
\usepackage[swedish]{babel}
\usepackage{lmodern}
\usepackage[autostyle]{csquotes}
\usepackage{enumitem}
\usepackage{booktabs}
\usepackage{longtable}
\usepackage{array}
\usepackage[backend=biber,style=authoryear,maxcitenames=2,
  uniquelist=false,useprefix=true]{biblatex}
\addbibresource{litteratur.bib}
\DefineBibliographyStrings{swedish}{andothers={m\adddot fl\adddot}}
\usepackage[hidelinks]{hyperref}
\usepackage{xurl}
\setlist{itemsep=2pt}
\usepackage{cleveref}
\usepackage{didactic}

\title{Delutvärdering DD1317 HT26 (prgi26)}
\author{Sammanställning, kommentarer och åtgärder}
\date{7 oktober 2026}

\begin{document}
\maketitle

\section{Sammanfattning}

Hittills har 18 av kursens 179 registrerade studenter svarat (cirka
10\,\%). Med så få svar, och med en trolig överrepresentation av de
missnöjda, är resultaten signaler snarare än andelar. Flera önskemål drar
dessutom åt olika håll (mer interaktion mot mer
\enquote{straight-forward} föreläsningar, fler exempel mot färre).

\begin{itemize}
\item \textbf{Svårighet och belastning:} 9 tycker att kursen är balanserad,
  7 svår, 1 mycket svår och 1 lätt. 11 tycker att programmeringskursen har
  lagom belastning, 2 att den är mycket hög. 7 skriver att de måste lägga
  mindre tid på programmeringen för att hinna med andra kurser.
\item \textbf{Föredraget material:} flest vill göra övningsuppgifter på
  egen hand (8) och se förinspelade videor (7); liveföreläsningar,
  föreläsningsanteckningar och material utanför kursen har 5 var,
  liveövningar 4.
\item \textbf{Starkaste temat:} föreläsningarna upplevs som ostrukturerade
  av nybörjare. De ber om en översikt av vad som ska läras, tydliga
  övergångar mellan ämnen och konstruktionens form innan den används.
\item \textbf{Canvas:} rummet upplevs som rörigt, och det är oklart vad som
  är obligatoriskt och vilka deadlines som gäller.
\item \textbf{Förkunskaper:} 56\,\% av de 156 som besvarat
  bakgrundsenkäten har aldrig programmerat
  (avsnitt~\ref{sec:forkunskaper}).
\item \textbf{Det jag inte förstår:} tre svar, om att föreläsningarna
  inte kodar live, om \enquote{mindre live och digitalt} och om att
  upplägget skulle vara till för att läraren ska arbeta mindre
  (avsnitt~\ref{sec:forstar-inte}).
\item \textbf{Önskemål om stöd för eget arbete:} en kort lista per modul
  med vad man ska kunna, övningsuppgifter med lösningar och en kortversion
  av anteckningarna.
\end{itemize}

\section{Studenternas förkunskaper}\label{sec:forkunskaper}

Bakgrundsenkäten (\emph{Diagnostic quiz: your programming background}),
som alla studenter besvarar innan Python-delen öppnas, ger en bild av
vilka som läser kursen. Siffrorna gäller alla 156 som svarat; i
forskningen används bara svaren från de 123 som samtyckt.
Översikt i \cref{tab:forkunskaper}.

\begin{table}\small
  \begin{fullwidth}
\begin{minipage}[t]{0.5\linewidth}\centering
\begin{tabular}{lr}
\toprule
Hur lärde du dig programmera? & Antal \\
\midrule
Jag har inte programmerat före kursen & 88 (56\,\%) \\
I en kurs & 27 \\
På egen hand & 22 \\
Både och & 16 \\
Inget svar & 3 \\
\bottomrule
\end{tabular}
\end{minipage}%
\begin{minipage}[t]{0.48\linewidth}\centering
\begin{tabular}{lr}
\toprule
Kan du skriva, köra och rätta & \\
ett litet program på egen hand? & Antal \\
\midrule
Inte alls säker & 82 \\
Lite säker & 34 \\
Ganska säker & 15 \\
Rätt säker & 10 \\
Mycket säker & 13 \\
Inget svar & 2 \\
\bottomrule
\end{tabular}
\end{minipage}
\end{fullwidth}
    \caption{Förkunskaper i programmering, bakgrundsenkäten DD1317 HT26.}
    \label{tab:forkunskaper}
\end{table}

Drygt hälften har aldrig programmerat, och de 42\,\% som har det
överskattar troligen erfarenheten: fritextsvaren på samma enkät
innehåller sådant som \enquote{minimalt med python, effektiv arbetstid
<10 timmar}. Tre av fyra är inte alls eller bara lite säkra på att de
kan skriva och rätta ett litet program. Det förklarar varför
nybörjarperspektivet är så tydligt i delutvärderingen: kursen ska i
första hand fungera för den som börjar från noll, medan de erfarna (som
ofta inte går på undervisningen) behöver andra utmaningar.

\section{Teman i fritextsvaren}

Varje tema avslutas med min kommentar, om jag har någon.

\subsection{Föreläsningarnas struktur}

Det mest återkommande temat, från både dem som går och dem som inte går.
Ord som \enquote{luddigt}, \enquote{flummigt}, \enquote{spretigt},
\enquote{osammanhängande} och \enquote{långrandigt} förekommer. Konkreta
önskemål:
\begin{itemize}
\item en lista med huvudrubriker (\enquote{typ på tavlan}) över vad
  föreläsningen ska lära ut, och tydliga markeringar när den går över till
  ett nytt ämne;
\item formen först: vad konstruktionen heter, när den används, hur den
  skrivs (rader, indentering, kolon, parenteser), vad som brukar följa med
  den, och \emph{sedan} experimenterandet;
\item att det nya inte trängs ihop under sista halvtimmen; ett svar påpekar
  att veckans material ofta inte hinns med på passet, så att det krävs
  självstudier;
\item att läraren kodar live och förklarar varje steg.
\end{itemize}
Samtidigt får föreläsningarna beröm för att de gör tydligt vad som
förväntas och hur teorin används, och interaktionen med klassen beskrivs
som bra när den repeterar förra veckans material, men svår att hänga med
i när allt är nytt för den som aldrig programmerat.

\emph{Kommentar.} Det här ligger i spänning med materialets design
(pröva först, förklara sedan): det som är tänkt som eget upptäckande
upplevs av en del nybörjare som avsaknad av struktur. Rubrikerna på tavlan
och formen före experimenterandet gjorde jag i princip redan på
föreläsningen den 6 oktober; det fortsätter (avsnitt~\ref{sec:planerat}).

\subsection{Nybörjarperspektivet}

Flera svar menar att förkunskaper förutsätts: \enquote{Det antas att
studenterna redan kan allt. Kanske vara lite mer övertydlig även om det
känns löjligt.} Man ber om grundligare repetition (\enquote{inte bara
\textquote{som ni känner till kan man använda valrossoperatorn}}), om att få veta
\enquote{lite mer exakt vad funktioner innebär på djupet}, och om
fler och kortare exempel på enskilda saker (\enquote{ex listor}), eftersom
exempel som bygger vidare på varandra blir svåra att följa om man missat
ett tillfälle. Bakgrundsenkäten (avsnitt~\ref{sec:forkunskaper}) visar
att detta gäller strax över hälften i kursen.

\subsection{Eget arbete och interaktion}

Flera vill öva mer själva: en kort uppgift med en tidsgräns, alla löser
den i salen, sedan genomgång, \enquote{kanske 50/50 det och att du visar
hur man gör som nu}. Interaktiva moment där man bygger riktiga program
får beröm, liksom att en assistent förklarar samma sak på ett annat sätt.

\emph{Kommentar.} Min erfarenhet talar emot tidsatta enskilda uppgifter i
salen: när jag på en övning bad studenterna arbeta med ett problem gick
alla i stället därifrån för att arbeta med laborationen. Folk går när de
ska jobba själva. Korta frågor i Mentipy under passet kan ge samma
aktivitet utan att salen töms (avsnitt~\ref{sec:planerat}). Men jag ska
prova igen med andra uppgifter.

\subsection{Laborationerna och vad man ska kunna}

Ett svar vill ha mer fokus på \enquote{hur man ska klara labbarna}, och
upplever att undervisningen ibland blir \enquote{nischad}. Labbarna
beskrivs som tydliga.

\emph{Kommentar.} Laborationerna är ett minimum, inte kursens innehåll.
De är medvetet enkla, och den som bara lär sig det som behövs för dem lär
sig för lite; om de skulle täcka allt man behöver kunna skulle de behöva
vara mycket större. Det behöver sägas till studenterna
(avsnitt~\ref{sec:planerat}).

\subsection{Översikt: vad ska jag kunna?}

Önskemål om en kort lista per modul med vad man förväntas kunna,
uppgifter som övar just det, exempel på bra lösningar och en
\enquote{nedkokad} version av anteckningarna att repetera från. Ett svar
vill ha uppmaningar och övningar att göra medan man läser. Anteckningarna
i övrigt får mest omdömen som \enquote{bra} och \enquote{ok}.

\emph{Kommentar.} Två av önskemålen finns redan men hittas inte:
anteckningarnas sammandrag anger vad man ska kunna efter varje
föreläsning, och övningsanteckningarna är just uppgifter med
lösningsförslag. Det verkar behöva pekas ut tydligare.

\subsection{Videorna}\label{sec:videorna}

Ett svar konstaterar att föreläsningen \enquote{känns som att det är
samma som på videorna}. Ett annat berömmer videorna som \enquote{riktigt
bra och tydligt hur man kan applicera dem till sin labb}, och föredrar
dem framför föreläsningen eftersom man \enquote{själv [får] tid att
processera infot}.

\emph{Kommentar.} Ja, och det är avsiktligt: i år ändrades upplägget så
att inspelningen är föreläsningen, för den som inte kan vara på plats.
Tidigare år efterfrågade studenterna mer traditionella föreläsningar
framför ett flippat klassrum, trots att flippat klassrum i genomsnitt ger
något bättre studieresultat än föreläsningsbaserad undervisning
\parencite{LagSaele2019,vanAlten2019}.\footnote{Hur stor effekten är, vad
den beror på och vad som talar emot den redovisas i
\hyperref[app:flippat]{bilaga~A}, med litteratursökningen bakom.}
Effekten är liten, omkring en femtedel till en tredjedel av en
standardavvikelse, varierar kraftigt mellan kurser och är störst när
förberedelserna följs upp med quiz \parencite{vanAlten2019}; när även den
traditionella undervisningen innehåller aktiva moment krymper skillnaden
eller försvinner \parencite{Jensen2015,Kapur2022}. För programmering är
underlaget tunt: den enda metaanalysen som sökningen fann bygger på 14 studier med mycket
stor spridning mellan dem \parencite{AlmassriZaharudin2023}.

Att studenter föredrar föreläsningar trots att mer aktiv undervisning ger
bättre resultat är belagt: i ett randomiserat försök i en inledande
fysikkurs presterade studenterna bättre efter aktivt undervisade lektioner
men upplevde att de lärt sig mindre och föredrog föreläsningarna
\parencite{Deslauriers2019}. Mönstret är dock inte allmängiltigt; i ett
liknande försök inom läkarutbildningen upplevde studenterna tvärtom att de
lärde sig mer i de aktiva passen \parencite{Boedeker2025}. Årets upplägg,
med föreläsning följd av övning och laboration, har redan aktiva moment,
så skillnaden mot ett flippat klassrum bör vara liten.

\subsection{Canvas och kursens administration}

Rummet upplevs som rörigt; det är svårt att se vad som måste lämnas in
eftersom många moment har poäng, och deadlines och datum för inlämning
och redovisning upplevs som oklara. Ett svar ber om att frågor som ställs
på Canvas besvaras, ett annat om att tidslåsen på uppgifter tas bort.

\emph{Kommentar.} Frågor ska besvaras, och assistenterna ska göra det på
passen också.

\subsection{Komma igång}

Terminalen, installationen av en textredigerare, learnlog samt Git och
GitHub nämns som det krångligaste, och Git och GitHub upplevs fortfarande
som oklart. En student önskar en enkel checklista per operativsystem
(\enquote{Har du MAC -- gör denna to do lista}).

\subsection{Varför man inte går på undervisningen}

12 av 18 förklarar varför de inte går: tidigare erfarenhet, hög belastning
i andra kurser, att föreläsningarna upplevs som otydliga eller långa, och
att videor, YouTube eller AI ger mer per tidsenhet.

\section{Det jag inte förstår}\label{sec:forstar-inte}

Tre svar förstår jag inte, eller känner inte igen mig i. De behöver
förstås innan de kan mötas; ett sätt är att fråga efter förtydligande i en
uppföljning, eller uppmana studenterna att ta upp dem med kursnämnden.

\begin{enumerate}
\item \enquote{De är osammanhängande och oklara. Jag skulle vilja se när
  undervisaren kodar live och förklarar varje funktion och steg.}
  (Bilaga~B.2.3, svar~7.)

  Det är väl vad vi gör på föreläsningarna? Möjliga läsningar: studenten
  har varit på få pass (\enquote{ett fåtal}) och träffat dem där
  diskussionen kring en pröva-först-fråga tog tiden; eller det som visas
  är färdig kod på bilder snarare än kod som skrivs i stunden; eller
  koden skrivs live men utan att varje steg sägs ut.
\item \enquote{Mindre live och digitalt samt ett mindre rörigt
  canvasrum.} (Bilaga~B.2.4, svar~11.)

  Oklart vad som ska bli mindre. Möjliga läsningar: färre parallella
  kanaler (sal, Zoom, inspelning, anteckningar), så att det är tydligt
  var undervisningen sker; eller mindre av det samtidiga upplägget i sal
  och på Zoom; eller mindre av allt utom föreläsning plus övning. Resten
  av svaret ber om \enquote{tydliga och straight-forward föreläsningar
  med efterföljande övningar där man får se någon annan skriva kod}, vilket är i princip vad vi gör.
\item \enquote{Förstår att du vill köra mycket digitalt för att kunna
  jobba mindre men om en oerfaren kodare faktiskt ska bli bra på
  programmering tror jag att det behövs mer fysisk och tydlig
  undervisning} (Bilaga~B.2.6, svar~5), och den liknande tolkningen i
  ett annat svar.

  Uppfattningen att upplägget finns till för att jag ska arbeta mindre
  känner jag inte igen. Föreläsningar och övningar ges på plats i sal,
  och inspelningarna och Zoom är till för den som inte kan komma. Men
  kursen förklarar inte själv varför den är upplagd som den är
  (inspelningar, Zoom, kamratgranskning), så studenterna fyller tomrummet
  med en egen förklaring; det kan mötas med en kort motivering.
\end{enumerate}

\section{Åtgärder}

\subsection{Genomförda och påbörjade}

\begin{itemize}
\item Alla laborationer har deadline fredag kl.~19 den vecka ämnet tas
  upp (laboration~1--3 får datum som redan passerat).
\item Alla obligatoriska moment heter nu \enquote{\dots{}
  (obligatorisk)}: laborationerna, projektets moment, datorprovet,
  Excel-laborationerna och de diagnostiska proven.
\item Uppgiftsgruppen KAL1 innehåller bara Excel-laborationerna, och nya
  uppgifter hamnar inte längre där.
\item Formuleringar om gruppinlämning är borttagna: var och en lämnar in
  sin egen lösning.
\item Föreläsningen den 6 oktober började med huvudrubrikerna och visade
  konstruktionernas form innan de användes.
\item Sådant som redan görs och får beröm: interaktionen när
  föreläsningar och övningar repeterar förra veckans material, genomgångar där vi bygger riktiga program, en assistent som
  förklarar samma sak på ett annat sätt, korta exempel på enskilda
  konstruktioner, och svar på frågor i Canvas.
\end{itemize}

\subsection{Planerade}\label{sec:planerat}

\begin{enumerate}
\item \textbf{Översikt i början och slutet av varje föreläsning}:
  huvudrubrikerna och \enquote{det här ska du kunna efter föreläsningen}
  från lärandemålen, gärna som Mentipy-frågor i början och slutet, så att
  studenterna själva ser vad de redan kan. Signalera övergångar muntligt.
\item \textbf{Syntaxruta efter varje avslöjande}: konstruktionens form,
  när den används och vad som brukar följa med den, efter
  pröva-först-momentet; veckorna 44--46 först.
\item \textbf{Peka ut det som redan finns}: på varje veckosida en
  hänvisning till anteckningarnas sammandrag (vad man ska kunna) och till
  övningsanteckningarna (uppgifter med lösningsförslag).
\item \textbf{Aktivitet på passen}: Mentipy-frågor, och ett nytt försök
  med andra slags uppgifter att arbeta med själv, eftersom studenterna
  gick när de fick ett problem att lösa på egen hand.
\item \textbf{Ett meddelande till studenterna} om vad vi hört och vad som
  ändras, med en kort förklaring av upplägget: inspelningarna är
  föreläsningen för den som inte kan komma, kamratgranskningen är en del
  av det man lär sig, och laborationerna är ett minimum, inte allt man
  ska kunna. Meddelandet kan också be om förtydligande av svaren i
  avsnitt~\ref{sec:forstar-inte}, eller hänvisa till kursnämnden, och
  be fler svara.
\item \textbf{Sammanfattning skriver studenterna själva.} Önskemålet om en
  kortversion möts med att anteckningarnas sammandrag redan anger vad man
  ska kunna, och med en uppmaning att skriva en egen sammanfattning inför
  datorprovet, som en del av studierna.
\item \textbf{Assistenterna svarar på frågor} på Canvas och på passen.
\item \textbf{Checklista för att komma igång} per operativsystem, med Git
  och GitHub.
\end{enumerate}

\subsection{Till nästa år}

Ett omstrukturerat Canvasrum (färre moment med poäng, en plats för
deadlines), en enklare start med learnlog, och en översyn av hur mycket
pröva-först nybörjarna får, med årets data som underlag.

\printbibliography[heading=bibintoc,title={Referenser}]

\appendix
% Bilaga A: the literature review behind §3.6 (Videorna).
%
% Author's note (verification and reproduction):
%   scholar session : flipped-classroom-learning
%   exports         : /tmp/claude-1000/research-flipped/hits.csv (classified,
%                     845 unique records), hits-table.tex (bearing-only
%                     longtable, not \input here; see tab:flippat-kallor)
%   searched        : 2026-10-07; providers openalex, wos, scopus, ieee
%                     (scopus standard view, no abstracts);
%                     s2: key rejected (HTTP 403); dblp: every request
%                     answered with an Anubis bot-challenge page (non-JSON)
%   queries         : S1 S2 S3 C1 C2 C3 P1 P2 + P1w (P1 with truncation, on
%                     ieee/wos/scopus); verbatim in the header of
%                     litteratur.bib.  In the text C1-C3 are M1-M3.
%                     The first round (same queries) was rerun: its WoS
%                     cells used TS=(FC) AND ... so TS covered only the
%                     first group (400 errors), and OpenAlex rejected
%                     wildcards inside phrases.  Coverage of the redone
%                     search checked by normalised title against the first
%                     round: all first-round items present after adding
%                     P1w (2026-10-07).
%   limit           : -l 50 per query and provider
%   classification  : scholar llm context + llm classify --no-examples,
%                     model openai-codex/gpt-5.6-sol, tags supports-claim /
%                     refutes-claim / qualifies-claim / adjacent-subtopic /
%                     off-topic-false-hit; every refutes-claim row and every
%                     cited source read; one row below confidence 0.6.
%                     Cited sources recorded with scholar sessions decide
%                     --keep and theme tags metaanalys, metodkritik,
%                     programmering, preferens.
%   full texts      : scholar PDF cache (Låg & Sæle via ERIC EJ1229661,
%                     van Alten via Utrecht repository, Hew et al. 2021
%                     via figshare, Kapur, Almassri) and Europe PMC XML
%                     (Deslauriers PMC6765278, Hew & Lo PMC5855972, Jensen
%                     PMC4353080, Boedeker PMC11933506).
%   provenance      : litteratur.bib, a CLAIM/FOUND-VIA/PICKED/QUOTE/
%                     VERIFIED/COUNTER/DATE block above every entry.

\clearpage
\phantomsection\label{app:flippat}
\section*{Bilaga A: Ger ett flippat klassrum bättre lärande än
  föreläsningar?}
\addcontentsline{toc}{section}{Bilaga A: Ger ett flippat klassrum bättre
  lärande än föreläsningar?}
\renewcommand{\thesection}{A.\arabic{section}}
\renewcommand{\theHsection}{bilagaA.\arabic{section}}
\renewcommand{\thetable}{A.\arabic{table}}
\renewcommand{\theHtable}{bilagaA.\arabic{table}}
\setcounter{section}{0}
\setcounter{table}{0}

\emph{Författaren har ännu inte granskat resultaten i den här bilagan i sin
helhet.}

\section{Frågan}

Avsnitt~\ref{sec:videorna} säger att ett flippat klassrum i genomsnitt ger
något bättre studieresultat än föreläsningsbaserad undervisning
\parencite{LagSaele2019,vanAlten2019}, att effekten är liten och att den
krymper när även den traditionella undervisningen har aktiva moment. I ett
flippat klassrum tar studenterna del av genomgången, oftast som video, före
passet, och passet används till att arbeta med stoffet. Bilagan prövar
påståendet genom att ställa frågan om själva sakförhållandet: ger ett
flippat klassrum bättre lärande än föreläsningsbaserad undervisning, i
högre utbildning i allmänhet och i inledande programmering i synnerhet?
Till den hör tre följdfrågor: hur stor och hur säker effekten är, vilka
drag i upplägget som bär den, och hur studenternas nöjdhet och upplevda
lärande förhåller sig till det faktiska lärandet. Påståendet som prövas är
\enquote{i genomsnitt något bättre}, inte \enquote{alltid bättre}.

\section{Metod}

Sökningen gjordes den 7 oktober 2026 med verktyget \texttt{scholar}, som
skickar samma fråga till flera bibliografiska databaser och sparar frågor
och träffar. Fyra databaser svarade: OpenAlex, Web of Science, Scopus och
IEEE Xplore; Scopus lämnade bara grunduppgifter, utan sammanfattningar.
Två databaser föll bort: DBLP svarade på varje fråga med en sida som
kontrollerar att besökaren inte är en robot, och Semantic Scholar avvisade
verktygets nyckel. Deras kolumner saknas därför i
tabell~\ref{tab:flippat-fragor}; de betyder inte noll träffar.

Frågorna byggs av samma begreppsblock. Flippat klassrum (FC) är
\enquote{flipped classroom} OR \enquote{inverted classroom} OR
\enquote{flipped learning}; utfallen är achievement, \enquote{learning
outcomes}, performance och grades; översikterna är meta-analysis,
\enquote{systematic review} och \enquote{literature review}; ämnet är
programming, \enquote{computer science} och computing. Stödfrågorna (S)
söker översikter och jämförelser. Motfrågorna (M) återanvänder blocken och
lägger till det som skulle försvaga påståendet: publiceringsbias,
heterogenitet, moderatorer, begränsningar, \enquote{no significant
difference}, negativa resultat och motstånd. Två frågor (P) gäller
studenternas preferenser och upplevda lärande. Web of Science fick varje
fråga som \texttt{TS=(\dots)}, Scopus som \texttt{TITLE-ABS-KEY(\dots)},
OpenAlex och IEEE Xplore som boolesk fritext. Fråga P1 kördes med
trunkering (\texttt{lecture*}, \texttt{preference*}) i de databaser som
tillåter det och med utskrivna böjningsformer i OpenAlex. Varje fråga gav
högst 50 träffar per databas, i databasens relevansordning.

\begin{table}[htbp]
  \centering\footnotesize
  \caption{Sökfrågor och antal träffar per databas, 7 oktober 2026, med
    verktyget \texttt{scholar}. S = stödfråga, M = motfråga, P = fråga om
    preferens och upplevt lärande. FC står för (\enquote{flipped
    classroom} OR \enquote{inverted classroom} OR \enquote{flipped
    learning}); (s) betyder att både singular och plural söktes.
    Databaser: OpenAlex (OA; inte \foreignlanguage{english}{open access}),
    Web of Science (WoS), Scopus och IEEE Xplore (IEEE). Högst 50 träffar
    per databas och fråga, så 50 betyder \enquote{minst 50}. DBLP och
    Semantic Scholar svarade inte och saknas.}
  \label{tab:flippat-fragor}
  \begin{tabular}{@{}l>{\raggedright\arraybackslash}p{0.62\linewidth}rrrr@{}}
    \toprule
    & Fråga & OA & WoS & Scopus & IEEE \\
    \midrule
    S1 & FC AND (meta-analysis OR \enquote{systematic review}) AND
      (achievement OR \enquote{learning outcomes} OR performance)
      & 50 & 50 & 50 & 18 \\
    S2 & FC AND (programming OR \enquote{computer science} OR computing)
      AND (achievement OR \enquote{learning outcomes} OR performance OR
      grades) & 50 & 50 & 50 & 50 \\
    S3 & FC AND (programming OR \enquote{computer science} OR computing)
      AND (meta-analysis OR \enquote{systematic review} OR
      \enquote{literature review}) & 50 & 50 & 47 & 28 \\
    M1 & FC AND (meta-analysis OR \enquote{systematic review}) AND
      (\enquote{publication bias} OR heterogeneity OR \enquote{no
      significant difference} OR moderator(s) OR limitations)
      & 50 & 50 & 50 & 6 \\
    M2 & FC AND (programming OR \enquote{computer science} OR computing)
      AND (\enquote{no significant difference} OR \enquote{no difference}
      OR challenges OR resistance OR negative) & 50 & 50 & 50 & 50 \\
    M3 & S3 AND (\enquote{no significant difference} OR heterogeneity OR
      limitations OR \enquote{mixed results}) & 50 & 8 & 1 & 1 \\
    P1 & (FC OR \enquote{active learning}) AND lecture(s) AND
      (\enquote{student perception(s)} OR satisfaction OR preference(s) OR
      resistance OR \enquote{feeling of learning}) & 50 & 50 & 50 & 50 \\
    P2 & (FC OR \enquote{active learning}) AND (\enquote{perceived
      learning} OR \enquote{feeling of learning} OR \enquote{self-reported
      learning}) AND (lecture(s) OR passive) & 50 & 33 & 40 & 6 \\
    \bottomrule
  \end{tabular}
\end{table}

Sökningen gav 845 unika träffar. En språkmodell klassade var och en mot en
beskrivning av frågan, som stöder påståendet, nyanserar det, motsäger det,
angränsande ämne eller felträff, med en angiven säkerhet; en träff kunde
få både \enquote{stöder} och \enquote{nyanserar}. Varje träff som modellen
klassade som motsägande lästes på sammanfattningsnivå, liksom varje källa
som bilagan bygger på; bara en träff hade en säkerhet under 0,6.

Källorna valdes i fyra lager: först de metaanalyser som jämför flippat
klassrum med föreläsningsbaserad undervisning över flera ämnen, sedan
granskningar av metaanalysernas metod och kontrollerade jämförelser som
talar emot eller nyanserar, sedan det som rör programmering, datavetenskap
eller matematik, och sist studier av preferens och upplevt lärande. Varje
källa lästes i fulltext där en öppen version fanns, hos förlaget, i ett
lärosätes arkiv eller i Europe PMC; annars lästes sammanfattningen.
Tabell~\ref{tab:flippat-kallor} anger vilket för varje källa. Siffror ur
en fulltext anges med sidnummer; fyra fulltexter fanns bara utan
sidnumrering, och för dem anges sidan bara där den framgår. Siffror som
bara kommer från en sammanfattning sägs komma därifrån.

\section{Resultat}

\subsection{Effekten i högre utbildning}

Den bredaste metaanalysen, av \textcite{LagSaele2019}, jämför flippat
klassrum med traditionell föreläsningsbaserad undervisning i 272 studier
från alla nivåer och ämnen. Den sammanvägda effekten på lärande är
\(g = 0{,}35\), med 95\,\% konfidensintervall 0,31--0,40
\parencite[1]{LagSaele2019}. Med bara de studier som har tillräcklig
statistisk styrka blir den 0,24, och korrigerad för saknade små studier
med metoden \foreignlanguage{english}{trim and fill} 0,17; författarna
sammanfattar: \enquote{the true weighted mean effect is somewhere around
one fifth of a standard deviation} \parencite[9]{LagSaele2019}. Andelen
godkända är 77,6\,\% mot 75,9\,\% \parencite[7]{LagSaele2019}.
\textcite[s.~8 och tabell~3, s.~10]{vanAlten2019} finner \(g = 0{,}36\)
(0,28--0,44) över 114 studier i gymnasium och högre utbildning.

Övriga metaanalyser pekar åt samma håll. Enligt sammanfattningarna finner
\textcite{Cheng2019} \(g = 0{,}19\) över 55 publikationer,
\textcite{Strelan2020} \(g = 0{,}50\) över 198 studier, och
\textcite{Bredow2021}, med 317 studier i högre utbildning, signifikanta
fördelar för flippat klassrum i sju av åtta utfall, med \(g\) mellan 0,20
och 0,53. Inom vårdutbildningar finner \textcite{HewLo2018} en
standardiserad medelvärdesskillnad på 0,33 (0,21--0,46). En metaanalys av
15 metaanalyser i högre utbildning ger \(g = 0{,}45\) (0,37--0,53), och
0,37 efter korrigering för publiceringsbias
\parencite[140--141]{Hew2021second}; dess tabell över de ingående
analyserna bekräftar Chengs och Strelans siffror för högre utbildning,
0,193 och 0,48 \parencite[136]{Hew2021second}. En liten metaanalys inom
farmaci, med fem studier, finner däremot ingen signifikant skillnad i
tentamensresultat och talar om \enquote{minimal gains}
\parencite{Gillette2018}.

\subsection{Spridning, publiceringsbias och metodkvalitet}

Spridningen mellan studierna är mycket stor. Andelen av variationen som
inte beror på slumpen, \(I^2\), är 83\,\% hos
\textcite[7]{LagSaele2019}, 88\,\% hos \textcite[tabell~3,
s.~10]{vanAlten2019} och 91\,\% i den metaanalys som
\textcite[7]{Kapur2022} gör. Kapur m.fl.\ samlar också 46 metaanalyser vars sammanvägda effekter
ligger mellan 0,19 och 2,29 \parencite[2]{Kapur2022}. Publiceringsbias är
belagd: små studier drar upp effekten hos \textcite[9]{LagSaele2019},
och i
metaanalysen av metaanalyser ger både rangkorrelation och Eggers test
tecken på bias \parencite[141]{Hew2021second}. En granskning av 19
metaanalyser finner enligt sammanfattningen att bland annat sökning i grå
litteratur, bedömning av risk för bias, kontroll för studenternas
förkunskaper och för att samma lärare undervisar båda grupperna
\enquote{were inadequately conducted or reported}
\parencite{Hew2021methodology}.

\subsection{Vad effekten beror på}

Två drag i upplägget har stöd. Quiz på förberedelserna höjer effekten med
\(g = 0{,}19\), och en kortad lektionstid sänker den med 0,26
\parencite[9]{vanAlten2019}; \textcite{HewLo2018} finner också större
effekt när passet inleds med quiz. Hos \textcite[7]{LagSaele2019} är
skillnaden för prov på förberedelserna bara marginell (0,40 mot 0,31,
\(p = 0{,}051\)), och sociala aktiviteter i passet gör ingen påvisbar
skillnad.

Om det är det aktiva arbetet som bär effekten är omstritt. Enligt
sammanfattningen ser \textcite{Strelan2020} möjligheten till strukturerat
aktivt lärande som den främsta orsaken. \textcite{Kapur2022} kodar i
stället vad som faktiskt gjordes i 173 studier och finner att aktivt
lärande ofta saknades i de flippade passen och inte tillförde något när
det fanns; fördelen tycks komma av mer exponering och tid med stoffet.
De skriver: \enquote{As more active learning was incorporated in
traditional instruction, the effect of flipped learning over traditional
instruction tended to reduce, and in several cases, even reverse}
\parencite[12]{Kapur2022}. Deras redovisning har dock en motsägelse: på
s.~10 kallas skillnaden mellan pass med och utan aktivt lärande
(\(g = 0{,}42\) mot 0,33) obefintlig, samtidigt som testet anges till
\(t = 4{,}11\), \(p < 0{,}001\). \textcite{Jensen2015} jämförde i ett
kvasiexperiment två grupper i en biologikurs, 53 och 55 studenter med
samma lärare, där båda arbetade aktivt och bara ordningen skilde:
lärandet och nöjdheten blev desamma, och författarna drar slutsatsen att
vinsterna \enquote{are most likely a result of the active-learning style
of instruction rather than the order in which the instructor participated
in the learning process}.

\subsection{Programmering och datavetenskap}

Underlaget för programmering är tunt men pekar åt samma håll. Enligt
sammanfattningen är effekten för IT-ämnen den svagaste hos
\textcite{Strelan2020}, \(g = 0{,}30\) över 14 studier. Hos
\textcite[tabell~4, s.~11]{vanAlten2019} skiljer sig de formella
vetenskaperna, som artikeln inte närmare definierar, inte från
naturvetenskaperna (\(-0{,}035\), \(p = 0{,}698\)), och hos
\textcite[7]{LagSaele2019} är effekten för naturvetenskap, teknik och
matematik 0,32. Den enda metaanalysen för programmering som sökningen
fann, av \textcite[267, 276]{AlmassriZaharudin2023}, väger samman 14
studier till \(g = 0{,}56\) med \(I^2 = 90{,}1\,\%\). Den är svag:
konfidensintervallet anges olika i sammanfattningen och i texten
(0,33--0,79 mot 0,32--0,81, s.~275), och publiceringsbias bedöms bara med
trattdiagram och felsäkert \(N\). Två enskilda studier finner, enligt sammanfattningarna, ingen
fördel eller en nackdel: inga signifikanta skillnader i inledande
programmering över flera år \parencite{Tyler2018}, och mindre
lärandeökning med flippat klassrum i statistisk programmering, som
författarna förklarar med missuppfattningar som befästs under
självstudierna \parencite{Liu2026}. En tidig översikt av 32 artiklar om
flippat klassrum i datavetenskap efterlyser kontrollerade experiment
\parencite{Giannakos2014}, och en översikt för matematik finner stöd men
inkonsekventa resultat, beroende på upplägg och studentgrupp
\parencite{Jones2026}.

\subsection{Nöjdhet, preferens och upplevt lärande}

I genomsnitt påverkar upplägget inte nöjdheten. Hos
\textcite[9]{LagSaele2019} är effekten på nöjdhet 0,16 och efter
korrigering för saknade studier \(-0{,}03\), och hos
\textcite[tabell~3, s.~10]{vanAlten2019} 0,05, inte signifikant. Låg och
Sæle påpekar att flera
studier har effekter till föreläsningarnas fördel och att
\enquote{students may dislike an educational intervention even though it
improves learning} \parencite[12]{LagSaele2019}. I vårdutbildningarna
föredrog däremot fler studenter det flippade klassrummet
\parencite{HewLo2018}.

\textcite{Deslauriers2019} lottade studenter i en inledande fysikkurs
mellan aktivt undervisade lektioner och flytande föreläsningar med samma
material; det aktiva var arbete under passet, inte förberedelser före,
så upplägget är inte flippat. På provet låg de aktivt undervisade
studenterna nästan en halv standardavvikelse högre (0,46), men de
upplevde att de lärt sig mer än en halv standardavvikelse mindre (0,56)
och föredrog föreläsningarna: \enquote{students in the active classroom
learn more, but they feel like they learn less}
\parencite[19251]{Deslauriers2019}. \textcite{Boedeker2025} gjorde
ett liknande försök med 146 läkarstudenter och fann ett högre provresultat (0,27
standardavvikelser) men också ett högre upplevt lärande (0,56) i de
aktiva passen; de förklarar skillnaden mot Deslauriers m.fl.\ med
studenternas stadium eller innehållet. Två studier inom datavetenskap
mäter bara upplevelsen, enligt sammanfattningarna: i en flippad IT-kurs
var studenterna nöjda med upplägget vid kursens slut
\parencite{Elliott2014}, och i inledande programmering uppskattade
studenterna måttligt aktiva tekniker som livekodning mer än mycket aktiva
som kodning under passet \parencite{SoosaiRaj2018}.

\section{Diskussion och begränsningar}

Det som talar för påståendet är att alla breda metaanalyser finner en
positiv effekt, oberoende av varandra och med delvis olika studier. Det
som talar emot ett starkare påstående är att effekten är liten när den
korrigeras för publiceringsbias, att spridningen mellan studierna är så
stor att genomsnittet säger lite om en enskild kurs, att metaanalyserna
har kända metodbrister, och att en kontrollerad jämförelse och en
metaanalys som kodar vad som faktiskt gjordes i passen
\parencite{Jensen2015,Kapur2022} pekar på att fördelen krymper eller
försvinner när den icke-flippade undervisningen också är aktiv. Vad som
bär effekten är omstritt; quiz på förberedelserna och bibehållen
lärarledd tid är de drag som har stöd i mer än en studie.

För den här kursen betyder det att jämförelsen i de flesta studierna,
flippat klassrum mot föreläsning med hemuppgifter, inte är årets val.
Årets upplägg är föreläsning, live och inspelad, följd av övning och
laboration: en traditionell ordning med aktiva moment efter genomgången.
Mot ett sådant upplägg talar underlaget för en liten eller ingen skillnad,
och de drag som verkar bära effekten, uppföljda förberedelser och
lärarledd tid, kan finnas i båda. Studenternas önskan om
föreläsningar säger lite om lärandet: nöjdheten beror i genomsnitt inte på
upplägget, och upplevt och faktiskt lärande kan gå isär, även om de inte
alltid gör det.

Underlaget har begränsningar. Två av sex databaser föll bort, och ACM:s
konferenser om datavetenskaplig undervisning nås bara i den mån de andra
databaserna indexerar dem. Varje fråga gav högst 50 träffar per databas,
så källor längre ned i relevansordningen kan ha missats. Klassningen
gjordes av en språkmodell; bara de motsägande träffarna och de citerade
källorna lästes. Elva av de tjugo källorna lästes bara som sammanfattning
(tabell~\ref{tab:flippat-kallor}), bland dem fyra av metaanalyserna; för
Deslauriers m.fl.\ kunde bara en sida beläggas. De flesta studierna är
inte randomiserade, och många gäller en kurs första försök med flippat
klassrum \parencite[12]{LagSaele2019}. Tre metaanalyser av metaanalyser
från 2025 och 2026 hittades men kunde inte läsas.

\section{Slutsats}

Ja, i genomsnitt ger ett flippat klassrum något bättre lärande än
föreläsningsbaserad undervisning, men effekten är liten, omkring en
femtedel av en standardavvikelse när den korrigeras för publiceringsbias,
och den varierar så mycket mellan kurser att den inte kan tas för given i
en enskild kurs. För programmering pekar det tunna underlaget åt samma
håll. Effekten tycks bäras av drag som uppföljda förberedelser och bibehållen
lärarledd tid snarare än av själva ordningen; om aktivt arbete i passet
bidrar är omstritt. När den traditionella undervisningen har aktiva moment
är skillnaden liten eller ingen. Nöjdheten påverkas i genomsnitt inte; att studenter föredrar
föreläsningar trots att de lär sig mer av aktiv undervisning är visat i
ett randomiserat försök men inte i ett annat. Formuleringen i
avsnitt~\ref{sec:videorna}, \enquote{i genomsnitt något bättre}, håller;
en formulering utan förbehåll, \enquote{bättre för lärandet}, håller inte.

\section{Fortsatt arbete}

Den här sökningen lämnade följande ogjort. DBLP och Semantic Scholar bör
frågas när de svarar, och ACM:s digitala bibliotek bör sökas direkt.
Fulltexterna till \textcite{Strelan2020}, \textcite{Bredow2021},
\textcite{Cheng2019}, \textcite{Hew2021methodology}, \textcite{Liu2026},
\textcite{Giannakos2014}, \textcite{Gillette2018}, \textcite{Tyler2018},
\textcite{Jones2026}, \textcite{Elliott2014} och \textcite{SoosaiRaj2018}
behöver läsas, i första hand Strelan m.fl.\ (siffran för IT) och Bredow
m.fl.\ (vilket utfall som inte var signifikant). Tre metaanalyser av
metaanalyser från 2025 och 2026 hittades men har inte lästs alls, inte
ens som sammanfattning: Ma m.fl.\ (2025),\footnote{Ma, Fan och Ning,
\enquote{The effectiveness of the flipped classroom: A second-order
meta-analysis}, \emph{International Journal of Educational Research},
2025, \url{https://doi.org/10.1016/j.ijer.2025.102855}.} Uzun m.fl.\
(2026)\footnote{Uzun, Kaya, Ünal och Mutlu-Bayraktar, \enquote{Impact of
flipped classrooms on higher education student outcomes: A second-order
meta-analysis}, \emph{Studies in Educational Evaluation}, 2026,
\url{https://doi.org/10.1016/j.stueduc.2026.101653}.} och Wang m.fl.\
(2026).\footnote{Wang, Wang, Mo och Shen, \enquote{Flipped classroom and
learning outcomes: a second-order meta-analysis}, \emph{Educational
Technology Research and Development}, 2026,
\url{https://doi.org/10.1007/s11423-026-10686-z}.} De kan ändra
storleken på effekten men knappast riktningen.

Sökningen fann inga randomiserade jämförelser i inledande programmering,
och i synnerhet inga mellan ett flippat klassrum och en icke-flippad kurs
som också är aktiv, som den Jensen m.fl.\ gjorde i biologi. Den fann
inte heller någon studie som mäter både faktiskt och upplevt lärande hos
programmeringsnybörjare. Om en
sådan studie inte fann någon skillnad mellan flippat klassrum och
föreläsning med övningar skulle avsnitt~\ref{sec:videorna} kunna säga att
inget av uppläggen är bättre för lärandet; fann den en tydlig fördel för
det flippade klassrummet skulle den tala för att gå tillbaka till det.

\section{Källor och träffar}

Tabell~\ref{tab:flippat-kallor} listar de källor bilagan bygger på, med
hur de lästes, och tabell~\ref{tab:flippat-motsager} de övriga träffar som
språkmodellen klassade som motsägande. Hela den klassade träfflistan, 845
poster, är för lång för rapporten: av de 845 är 20 citerade, 96 klassade
som stödjande, 248 som nyanserande eller motsägande, 374 som angränsande
och 104 som felträffar, och 3 är dubbletter av citerade källor.

% Author's note: the curated tables below replace the bearing-only
% longtable (scholar sessions export flipped-classroom-learning -f table
% --bearing-only --lang sv --label flipped-traffar --theme ... ; 364 rows,
% about eight pages) at the team lead's option.  Counts above are from
% that export's caption.
\begingroup\footnotesize
\begin{longtable}{@{}>{\raggedright\arraybackslash}p{0.19\textwidth}%
    >{\raggedright\arraybackslash}p{0.16\textwidth}%
    >{\raggedright\arraybackslash}p{0.15\textwidth}%
    >{\raggedright\arraybackslash}p{0.38\textwidth}@{}}
  \caption{Källorna bilagan bygger på, efter roll. Läst anger om
    fulltexten eller bara sammanfattningen lästes.}
  \label{tab:flippat-kallor}\\
  \toprule
  Källa & Slag & Läst & Bidrag \\
  \midrule
  \endfirsthead
  \multicolumn{4}{@{}l}{\emph{Tabell~\thetable{} (forts.)}}\\
  \toprule
  Källa & Slag & Läst & Bidrag \\
  \midrule
  \endhead
  \bottomrule
  \endlastfoot
  \multicolumn{4}{@{}l}{\emph{Metaanalyser av flippat klassrum mot
    föreläsning}}\\
  \textcite{LagSaele2019} & metaanalys, 272 studier & fulltext &
    \(g = 0{,}35\); korrigerat 0,17--0,24; nöjdhet \(-0{,}03\) \\
  \textcite{vanAlten2019} & metaanalys, 114 studier & fulltext &
    \(g = 0{,}36\); quiz \(+0{,}19\), kortad tid \(-0{,}26\); nöjdhet
    0,05 \\
  \textcite{Hew2021second} & metaanalys av 15 metaanalyser & fulltext &
    \(g = 0{,}45\), 0,37 efter korrigering; tecken på publiceringsbias \\
  \textcite{HewLo2018} & metaanalys, 28 studier, vård & fulltext &
    0,33; större effekt med quiz i början av passet \\
  \textcite{Cheng2019} & metaanalys, 55 publikationer & sammanfattning &
    \(g = 0{,}19\) \\
  \textcite{Strelan2020} & metaanalys, 198 studier & sammanfattning &
    \(g = 0{,}50\); IT 0,30 \\
  \textcite{Bredow2021} & metaanalys, 317 studier & sammanfattning &
    fördel i sju av åtta utfall, \(g\) 0,20--0,53 \\
  \textcite{Gillette2018} & metaanalys, 5 studier, farmaci &
    sammanfattning & ingen signifikant skillnad \\
  \midrule
  \multicolumn{4}{@{}l}{\emph{Metodgranskning och kontrollerade
    jämförelser}}\\
  \textcite{Kapur2022} & 46 metaanalyser och ny metaanalys &
    fulltext & effekter 0,19--2,29; fördelen krymper när den
    traditionella undervisningen är aktiv \\
  \textcite{Hew2021methodology} & granskning av 19 metaanalyser &
    sammanfattning & brister i sökning, biasbedömning och kontroll \\
  \textcite{Jensen2015} & kvasiexperiment, biologi & fulltext &
    aktivt flippat = aktivt icke-flippat \\
  \midrule
  \multicolumn{4}{@{}l}{\emph{Programmering, datavetenskap och
    matematik}}\\
  \textcite{AlmassriZaharudin2023} & metaanalys, 14 studier & fulltext &
    \(g = 0{,}56\), \(I^2 = 90\,\%\); svag metod \\
  \textcite{Tyler2018} & kursjämförelse över år & sammanfattning &
    ingen signifikant skillnad \\
  \textcite{Liu2026} & jämförelse, statistisk programmering &
    sammanfattning & mindre lärandeökning med flippat \\
  \textcite{Giannakos2014} & översikt, 32 artiklar & sammanfattning &
    efterlyser kontrollerade experiment \\
  \textcite{Jones2026} & översikt, matematik & sammanfattning &
    stöd men inkonsekventa resultat \\
  \midrule
  \multicolumn{4}{@{}l}{\emph{Preferens och upplevt lärande}}\\
  \textcite{Deslauriers2019} & randomiserat försök, fysik & fulltext &
    lärde mer, upplevde mindre, föredrog föreläsning \\
  \textcite{Boedeker2025} & randomiserat försök, medicin & fulltext &
    lärde mer och upplevde mer \\
  \textcite{Elliott2014} & enkätstudie, IT & sammanfattning & nöjda vid
    kursens slut \\
  \textcite{SoosaiRaj2018} & enkätstudie, inledande programmering &
    sammanfattning & föredrog livekodning framför kodning under passet \\
\end{longtable}
\endgroup

\begingroup\footnotesize
\begin{longtable}{@{}>{\raggedright\arraybackslash}p{0.62\textwidth}%
    r>{\raggedright\arraybackslash}p{0.12\textwidth}r@{}}
  \caption{Övriga träffar som språkmodellen klassade som motsägande, med
    år, databas och modellens säkerhet. De lästes på
    sammanfattningsnivå men citeras inte; flera gäller andra upplägg än
    flippat klassrum, andra ämnen än programmering eller bara studenternas
    upplevelse.}
  \label{tab:flippat-motsager}\\
  \toprule
  Titel & År & Databas & Säkerhet \\
  \midrule
  \endfirsthead
  \multicolumn{4}{@{}l}{\emph{Tabell~\thetable{} (forts.)}}\\
  \toprule
  Titel & År & Databas & Säkerhet \\
  \midrule
  \endhead
  \bottomrule
  \endlastfoot
  East versus west: The descendants of confucianism vs. evidence-based learning mainland Chinese learners in pursuit of western-based education in Singapore & 2013 & IEEE & 0,920 \\
  Inverted Classroom by Topic - A Study in Mathematics for Electrical Engineering Students & 2014 & OA & 0,980 \\
  Switching to blend-Ed: Effects of replacing the textbook with the browser in an introductory computer programming course & 2016 & IEEE & 0,960 \\
  A Flipped Classroom Model For Algorithm In College & 2017 & OA & 0,980 \\
  Flipped classroom approaches lead to no improvement in learning outcomes or student perceptions & 2017 & OA & 0,990 \\
  The flipped classroom: A learning model to increase student engagement not academic achievement & 2017 & OA & 0,990 \\
  ATTITUDES FOR IMPROVED LEARNING \& KNOWLEDGE IN ARCHETYPAL ENGINEERING COURSES & 2019 & WoS & 0,960 \\
  Flipping General Chemistry in Small Classes: Students’ Perception and Success & 2019 & OA & 0,970 \\
  Introducing an Inverted Classroom Concept for an Undergraduate Course in Power Electronics – Implementation, Evaluation and Experiences & 2019 & IEEE & 0,970 \\
  Undergraduates’ satisfaction and perceptions of learning outcomes across teacher- and learner-focused pedagogies & 2019 & Scopus & 0,970 \\
  Probing the Inverted Classroom: A Controlled Study of Teaching and Learning Outcomes in Undergraduate Engineering and Mathematics & 2020 & OA & 0,870 \\
  The Effect of Motivation on Learners' Performance and Satisfaction under Flipped Strategy in Discrete Math & 2020 & IEEE & 0,970 \\
  An Examination of Student Preferences and Learning Outcomes in Flipped Classroom with Online Videos & 2022 & OA & 0,990 \\
  Dependence of learning outcomes in flipped and lecture classrooms on review questions: A randomized controlled trial and observational study & 2022 & Scopus & 0,990 \\
  The Effectiveness of Using MOOC with E-FLIC Approach in Remote Learning as New Normal of Education to Support Basic Computer Science Learning & 2025 & IEEE & 0,980 \\
  Contribution of Active Learning Methods in Large-Group Teaching of Medical Microbiology & 2026 & Scopus & 0,960 \\
  Simulation-based education in an emergency medicine clerkship in Qatar: impact on academic performance and student perceptions & 2026 & Scopus & 0,970 \\
  The use of voiceovers and review games in teaching immunology & 2026 & Scopus & 0,970 \\
\end{longtable}
\endgroup


\clearpage
\section*{Bilaga B: alla svar}
\addcontentsline{toc}{section}{Bilaga B: alla svar}
\renewcommand{\thesection}{B.\arabic{section}}
\renewcommand{\theHsection}{bilagaB.\arabic{section}}
\setcounter{section}{0}
\input{bilaga}

\end{document}
